\section{Pengenalan Environment Variable dan Konfigurasi Per Lingkungan}

\textbf{Environment variable} menyimpan konfigurasi yang berbeda antar lingkungan (dev, staging, prod) di luar kode sumber. Kredensial database, API key, mode debug, dan URL aplikasi disimpan dalam file \code{.env} yang \textbf{tidak boleh masuk ke version control} (tambahkan di \code{.gitignore}). Library \code{vlucas/phpdotenv} membaca file \code{.env} dan menyediakannya via \code{\$\_ENV} atau \code{getenv()}. Di server produksi, environment variable bisa diatur langsung di level server (Apache/Nginx) atau OS \cite{phpRightWay}.

\begin{webcode}[language=PHP, caption={File .env dan penggunaannya di PHP}]
# File: .env (JANGAN commit ke Git!)
APP_ENV=production
APP_DEBUG=false
DB_HOST=localhost
DB_NAME=toko_online
DB_USER=app_user
DB_PASS=p@ssw0rd_Kuat!123
DB_PORT=3306

# File: config.php
# <?php
# require 'vendor/autoload.php';
# $dotenv = Dotenv\Dotenv::createImmutable(__DIR__);
# $dotenv->load();
#
# $db_host = $_ENV['DB_HOST'];
# $db_name = $_ENV['DB_NAME'];
# $db_user = $_ENV['DB_USER'];
# $db_pass = $_ENV['DB_PASS'];
# $debug   = $_ENV['APP_DEBUG'] === 'true';

# File: .gitignore
# .env
# vendor/
# *.log
\end{webcode}

